\begin{frame}
\frametitle{Conclusion}
\textbf{What I showed:}
\begin{itemize}
  \item Time protection mechanisms \textit{can} be integrated into an existing mainstream OS
  \item Domain-level scheduler achieves constant-time scheduling in QEMU
  \item Trade-offs 
    \begin{itemize}
      \item \textbf{Complexity:} A significant amount of complexity has been added to the code. Future work onFreeRTOS would be more cumbersome with the extra code complexity.
      \item \textbf{Performance:} Given the implementation makes use of idle time, the total performance is most likely reduced significantly. Further, more memory accesses is needed for the domain Scheduler to function, so while asymptotically the same, it is a larger overhead.
    \end{itemize}
\end{itemize}

\textbf{Conclusion:} Full domain isolation is possible and functional in
general operating systems.

\end{frame}

\begin{frame}
\frametitle{Future work}
\begin{itemize}
  \item Evaluation on real hardware with \texttt{fence.t} support (or gem5)
  \item Review system calls: critical sections delay scheduling, only user-level protection
  \item Introduce a two level locking system.
    \begin{itemize}
      \item FreeRTOS has a notion of "critical section". These sections grab a
        global lock to ensure consistency of the priority-based scheduler. 
      \item Currently the domain level scheduler uses the same lock.
      \item \textbf{Idea:} Allow for a separate locking mechanism for
        cross-domain consistency. This would allow scheduling a new domain
        while holding the priority-based scheduler lock.
    \end{itemize}
  
  \item Remaining channels: interrupts, higher-level caches
\end{itemize}

\vspace{0.15cm}
%TODO: Rewrite this takeaway
Selective temporal partitioning might be more feasible in general operating systems.

\end{frame}
